\section*{Chapter \ref{ch:dv:variance-swaps-and-volatility-derivatives} --- Variance Swaps and Volatility Derivatives}

\begin{solution}{exo:dv:variance-swaps-and-volatility-derivatives:1}
$N_{\mathrm{var}}=50\,000/(2\times20)=1\,250$ per variance point; at 25\% the swap pays
$1\,250\times(625-400)=281\,250$.
\end{solution}

\begin{solution}{exo:dv:variance-swaps-and-volatility-derivatives:2}
The strip pays $-2\ln(S_T/F)$, which equals realised variance only after adding $2\int dS/S$: the gains of holding
$2/S_t$ in the index. Without the daily trade the position is a \omterm{def:dv:variance-swaps-and-volatility-derivatives:log}{log contract}, whose payoff depends on where the
index ends, not on the path.
\end{solution}

\begin{solution}{exo:dv:variance-swaps-and-volatility-derivatives:3}
$f'(S)=\frac2{S_0}\ln(S/S_0)$ and $f^{\prime\prime}(S)=\frac2{S_0S}$. Its strip weights are $2/(S_0K)$ against $2/K^2$ for the
\omterm{def:dv:variance-swaps-and-volatility-derivatives:vs}{variance swap}: relative to the \omterm{def:dv:variance-swaps-and-volatility-derivatives:vs}{variance swap} it weights each strike by $K/S_0$, less below the spot and more above.
On a downside-skewed index the expensive puts count for less, so the strike is lower.
\end{solution}

\begin{solution}{exo:dv:variance-swaps-and-volatility-derivatives:4}
Expected final variance $0.75\times18^2+0.25\times16^2=243+64=307$ against $22^2=484$: the value is $-177$ per unit of
\omterm{def:dv:variance-swaps-and-volatility-derivatives:varnot}{variance notional}.
\end{solution}

\begin{solution}{exo:dv:variance-swaps-and-volatility-derivatives:5}
$V=0.04$, $\Var=0.0004$: $0.2-0.0004/(8\times0.04^{1.5})=0.2-0.00625=19.4\%$.
\end{solution}

\begin{solution}{exo:dv:variance-swaps-and-volatility-derivatives:6}
The adjustment is proportional to the variance of realised variance divided by its level. At short expiries
variance has no time to move; at long expiries mean reversion ($\kappa=2.05$, a half-life of about four months)
averages its moves away. The variance of the average is largest when the expiry is a few half-lives long.
\end{solution}

\begin{solution}{exo:dv:variance-swaps-and-volatility-derivatives:7}
The full strip gives 16.32\%; the options worth at least 0.05 run from 87 to 106, and the number falls to 16.10\%,
0.22 point lower: the missing wings carried variance.
\end{solution}

\begin{solution}{exo:dv:variance-swaps-and-volatility-derivatives:8}
The hedge is exact only for continuous paths within the strip's range. Each jump leaves a loss of about
$|r|^3/3$ of variance, and a crash through the lowest strike leaves the payoff growing while the hedge stops. The
premium over the strip is the price of that jump and truncation risk (0.18 point on average under chapter 13's
Merton model, with a one-in-a-hundred loss of 1.9 points over three months), not profit.
\end{solution}

\begin{solution}{pb:dv:variance-swaps-and-volatility-derivatives:1}
\textbf{1.} 22.1\% and 19.2\%.
\textbf{2.} The puts: 72\% of the variance; strikes below 80 alone 24\%.
\textbf{3.} 2\,262 per variance point; 931\,000.
\textbf{4.} 1.6 to 1.8 points, depending on the spacing.
\textbf{5.} 20.8\%.
\textbf{6.} 21.4\%: the calibrated model misses the surface's wings, which the strip weights heavily.
\textbf{7.} 19.8\%.
\textbf{8.} 0.72, 1.62 and 1.10 points.
\textbf{9.} The volatility of variance $\eta$: 1.24 points at $\eta-0.1$, 2.01 at $\eta+0.1$.
\textbf{10.} Its payoff is the square root of realised variance, a non-linear function of the path's quadratic
variation; only functions of the terminal price are static strips, and variance is the one path functional that
reduces to one.
\textbf{11.} 17.1, 18.1 and 20.5.
\textbf{12.} $20.4-18.1=2.3$ points.
\textbf{13.} It slopes down, from 110\% at 80\% of the future to 94\% at 175\%.
\textbf{14.} With \omterm{def:dv:variance-swaps-and-volatility-derivatives:vs}{variance swaps}, dynamically: a \omterm{def:dv:variance-swaps-and-volatility-derivatives:volswap}{volatility swap}'s sensitivity to variance is about
$1/(2\sigma)$, so the desk holds \omterm{def:dv:variance-swaps-and-volatility-derivatives:vs}{variance swaps} in that ratio and adjusts it as volatility moves, a hedge whose
cost is the \omterm{def:dv:stochastic-volatility:sv}{volatility of volatility} the adjustment prices.
\textbf{15.} At least the jump correction, 0.18 point under chapter 13's jump model, plus a reserve for crashes
through the strip.
\textbf{16.} Volatility is the square root of variance, and the mean of a square root is below the square root of
the mean. The difference, here 1.6 points, is the price of the \omterm{def:dv:stochastic-volatility:sv}{volatility of volatility}.
\textbf{17.} The distribution of realised variance, above all the \omterm{def:dv:stochastic-volatility:sv}{volatility of volatility} $\eta$.
\textbf{18.} A reserve for the model: the spread of the adjustment across plausible $\eta$ (about 0.4 point either way for $\pm0.1$)
and the cost of re-hedging in a volatility spike.
\textbf{19.} 1.62 volatility points: 21.4\% against 19.8\% at one year.
\textbf{20.} Variance is replicated by a static strip and a delta; volatility needs a model.
\end{solution}

\begin{solution}{iq:dv:variance-swaps-and-volatility-derivatives:1}
Itô's formula gives $\int\sigma^2dt=2\int dS/S-2\ln(S_T/S_0)$. Hold a static strip of out-of-the-money options
with weights $2/K^2$ (the log payoff) and trade $2/S_t$ in the underlying daily; the strike is the strip's price
per unit of time.
\iqlookfor{the \omterm{def:dv:variance-swaps-and-volatility-derivatives:log}{log contract}, the $1/K^2$ weights and the dynamic leg.}
\end{solution}

\begin{solution}{iq:dv:variance-swaps-and-volatility-derivatives:2}
The strip weights options by $1/K^2$, so low strikes count more, and on an index the low strikes trade at higher
implied volatilities. The strike is an average of variance over the whole smile, tilted to the downside.
\iqlookfor{the weights and the skew.}
\end{solution}

\begin{solution}{iq:dv:variance-swaps-and-volatility-derivatives:3}
Below, by Jensen's inequality. To second order the gap is $\Var(\sigma_R^2)/(8V^{3/2})$: about 1.6 points at one year
under a \omterm{def:dv:stochastic-volatility:heston}{Heston model} calibrated to a typical index surface, more with a higher \omterm{def:dv:stochastic-volatility:sv}{volatility of volatility}.
\iqlookfor{the sign, the expansion, and that it needs a model.}
\end{solution}

\begin{solution}{iq:dv:variance-swaps-and-volatility-derivatives:4}
Find the forward from put--call parity, $K_0$ at or below it, take puts below and calls above (both averaged at
$K_0$), weight by $\Delta K/K^2$ with one-sided differences at the ends, subtract $(F/K_0-1)^2/T$, interpolate two
expiries to thirty days. Edge cases: the forward exactly on a strike; zero bids and where to cut the wings; uneven
strike spacing; stale quotes; expiry times in minutes.
\iqlookfor{the formula, the $K_0$ correction and the data cases.}
\end{solution}

\begin{solution}{iq:dv:variance-swaps-and-volatility-derivatives:5}
The swap's leg pays $r^2=0.04$ of variance for the day while the hedge earns $2(e^r-1-r)=0.0375$: a loss of about
$|r|^3/3$. If the gap goes below the strip's lowest strike the hedge also stops tracking. Losses scale with notional
and there is no dynamic fix; the protection is pricing the jump risk, caps, and buying deep puts.
\iqlookfor{the cubic term, the strip's range, and caps.}
\end{solution}

\begin{solution}{iq:dv:variance-swaps-and-volatility-derivatives:6}
The future is $\E[\mathrm{VIX}_T]$ and the forward variance-swap volatility is $\sqrt{\E[\mathrm{VIX}_T^2]}$; Jensen's
inequality makes the first smaller, by an amount set by the volatility of the index. To trade it: buy the \omterm{def:dv:rough-volatility-and-forward-variance-models:fwdvar}{forward
variance} (a calendar spread of \omterm{def:dv:variance-swaps-and-volatility-derivatives:vs}{variance swaps}) against a short future, which is long the volatility of the index,
and hedge its \omterm{def:dv:greeks-and-the-hedging-pnl:greeks}{vega} with index options; the position profits if realised volatility of the index exceeds what the
gap prices.
\iqlookfor{Jensen, and what the spread is exposed to.}
\end{solution}
